\section{The problem with a byte stream}
\label{sec:stream}

Two systems are connected by a serial link --- UART, RS-485, SPI, anything at all. One sends
bytes, the other receives them. That is all the link does.

Suppose the sending node wants to send the temperature 23 degrees and the humidity 45 per cent.
It sends two bytes:

\begin{console}
17 2D
\end{console}

The receiver gets two bytes. And now come the questions, one after the other:

\begin{itemize}
  \item \textbf{Where does the message begin?} If the receiver was connected in the middle of a
    transfer, \code{17} may well be the second byte of something else.
  \item \textbf{How long is it?} Are more bytes coming, or was that all?
  \item \textbf{What do they mean?} That \code{17} is the temperature and \code{2D} the humidity
    is an agreement that exists nowhere in the stream.
  \item \textbf{Did they arrive intact?} A disturbance that changed \code{17} into \code{16}
    looks exactly like a temperature of 22 degrees.
  \item \textbf{Who sent them, and to whom?} On a bus with more than two nodes that is not
    obvious.
\end{itemize}

None of the questions can be answered from the bytes themselves. The answers have to be added,
and what adds them is a \textbf{protocol}.

\subsection{Frame versus stream}

A \textbf{frame} is a delimited, self-describing unit in the stream. It says where it begins, how
long it is, what it contains, who it comes from, who it is for, and whether it arrived intact.

The stream is still nothing but bytes --- it always is. The difference is that the bytes are now
arranged according to an agreement both sides know about, and one that can be checked.

\begin{aside}
\note This is the only idea in the chapter, and the rest of the book is its consequences. Both the
protocol of your own in \kapref{ch:parsning} and CAN in \kapref{ch:can} solve the same five
questions. They differ in \emph{where} the solution sits: in your code, or in silicon.
\end{aside}
